\documentclass[11pt,a4paper]{article}

%\usepackage[utf8]{inputenc}
%\usepackage[T1]{fontenc}
\usepackage[polish]{babel}
\usepackage{geometry}
\usepackage{booktabs}
\usepackage{tabularray}
\usepackage{afterpage}
\usepackage{rotating}
\usepackage{longtable}
\usepackage{array}
\usepackage{xcolor}
\usepackage{hyperref}
\usepackage{titlesec}
\usepackage{parskip}
\usepackage{enumitem}
\usepackage{fancyhdr}
\usepackage{pdflscape}
\usepackage{makecell}
\usepackage[normalem]{ulem}% normalem: nie przejmuj \emph; \uline łamie się między wierszami

\usepackage{tikz}
\usetikzlibrary{decorations}
\usetikzlibrary{decorations.pathreplacing}
\usetikzlibrary{patterns}

\usepackage{tcolorbox}
\tcbuselibrary{breakable,skins}

\usepackage{caption}

\usepackage{mathtools}
\usepackage{etoolbox}
\usepackage{xpatch}
\usepackage{environ}% \RenewEnviron: przechwycenie treści dowodu do pominięcia (AZ_SKIP)

\usepackage[noend]{algpseudocode}
\usepackage{algorithm}
\usepackage{multirow}
\floatname{algorithm}{Algorytm}
\renewcommand{\listalgorithmname}{Spis algorytmów}

\usepackage{nicematrix}
\usepackage{cancel}
\usepackage{amsthm}
%\usepackage{amssymb}
\usepackage{amsmath}
\usepackage[warnings-off={mathtools-overbracket,mathtools-colon}]{unicode-math}
\setmathfont{XITS Math}
\mathtoolsset{mathic=true}

% Wariant ⩽/⩾ z linią równości pod tym samym kątem co < / > (jak \leqslant/\geqslant).
% Remap odłożony do \begin{document}, bo unicode-math finalizuje tablicę symboli
% dopiero przy starcie dokumentu i nadpisałby zwykły \let z preambuły.
\AtBeginDocument{%
	\let\leq\leqslant
	\let\geq\geqslant
	\allowdisplaybreaks
}

\geometry{margin=2.5cm, headheight=12.7pt}

% Ostatnia deska ratunku przy łamaniu akapitów: pozwala TeX-owi rozciągnąć
% trudny wiersz zamiast wypychać go poza margines (usuwa drobne Overfull \hbox
% w gęstym tekście z wtrąceniami matematycznymi).
\emergencystretch=2em
\hfuzz=0.6pt% toleruj przepełnienia poniżej ~0,2 mm (niewidoczne, a generują ostrzeżenia)

\definecolor{accent}{RGB}{30, 60, 100}
\definecolor{accentA}{RGB}{255, 187, 0}
\definecolor{accentB}{RGB}{255, 100, 0}
\definecolor{accentC}{RGB}{0, 140, 130}
\definecolor{placeholdercolor}{RGB}{160, 160, 160}
\definecolor{rowA}{RGB}{235, 242, 250}
\definecolor{dygresjagray}{RGB}{105, 105, 105}% treść dygresji — wyszarzona, materiał opcjonalny
\definecolor{obowiazkowybg}{RGB}{255, 252, 245}% tło dowodów obowiązkowych na egzaminie
\definecolor{obowiazkowyline}{RGB}{230, 200, 150}% nasycona linia na granicy podświetlenia
\definecolor{pomijalnybg}{RGB}{250, 250, 250}% tło dowodów pomijalnych przy powtórce
\definecolor{pomijalnyline}{RGB}{240, 240, 240}% szara linia na granicy podświetlenia
\definecolor{pomijalnytext}{RGB}{120, 120, 120}% przygaszony tekst dowodów pomijalnych

% Podświetlenie materiału obowiązkowego jako zamknięte pudełko w bloku tekstu
% (nie rozlewa się na marginesy). Delikatna ramka + akcentowy pasek z lewej dają
% czytelny sygnał „to obowiązkowe” bez agresywnego tła do krawędzi kartki.
% breakable: poprawnie łamie się między stronami (enhanced rysuje ładne górne/
% dolne krawędzie segmentów). Stranded rysunek może zostawić trochę tła u dołu
% łamanego segmentu, ale w zamkniętym pudełku jest to znacznie mniej rażące.
\newtcolorbox{obowiazkowybox}{%
	breakable, enhanced,
	colback=obowiazkowybg,
	colframe=obowiazkowyline,
	boxrule=0.5pt, % jednolita cienka ramka ze wszystkich stron
	arc=2pt, outer arc=2pt,
	left=10pt, right=10pt, top=8pt, bottom=8pt,
	boxsep=0pt,
	pad at break*=6pt,
}
% Usuń resztkowy odstęp pionowy na końcu ramki (np. \topsep zamykający proof),
% który potrafi przelać się na kolejną stronę jako pusty segment tcolorbox
% i zostawić blady pasek tła (colback rysuje się dla każdego segmentu).
\AtEndEnvironment{obowiazkowybox}{\par\removelastskip}
% Zarezerwuj miejsce PRZED otwarciem ramki. Inaczej \azNeedSpace twierdzenia
% wewnątrz świeżo otwartej ramki wymusza łamanie strony przy pustej zawartości,
% tworząc pusty pierwszy segment (blady pasek tła u góry/dołu strony).
\BeforeBeginEnvironment{obowiazkowybox}{\Needspace*{6\baselineskip}}

% ── WERSJA Z WYRÓŻNIENIAMI ───────────────────────────────────────
% Flaga \azhighlight włącza tło materiału obowiązkowego. Sterowana zmienną
% środowiskową AZ_HIGHLIGHT=1 (czytaną przez LuaTeX przy starcie) — ten sam
% plik źródłowy daje wersję zwykłą i wersję z wyróżnieniami.
\newbool{azhighlight}
\edef\azHighlightEnv{\directlua{tex.sprint(os.getenv("AZ_HIGHLIGHT") or "0")}}
\ifdefstring{\azHighlightEnv}{1}{\booltrue{azhighlight}}{}

% Flaga \azskip (AZ_SKIP=1): wycina nieobowiązkowe dowody w całości — wersja do
% szybkiej powtórki, w której zostają tylko dowody obowiązkowe.
\newbool{azskip}
\edef\azSkipEnv{\directlua{tex.sprint(os.getenv("AZ_SKIP") or "0")}}
\ifdefstring{\azSkipEnv}{1}{\booltrue{azskip}}{}

% Środowisko obowiazkowy: w wersji z wyróżnieniami otwiera ramkę
% obowiazkowybox, w zwykłej wersji jest przezroczyste (treść składana
% normalnie, bez tła). Dzięki temu twierdzenia obowiązkowe oznaczamy raz,
% a obie wersje generujemy z jednego źródła.
\newenvironment{obowiazkowy}
{\ifbool{azhighlight}{\begin{obowiazkowybox}}{}}
			{\ifbool{azhighlight}{\end{obowiazkowybox}}{}}

% Dowody nieobowiązkowe: w wersji z wyróżnieniami szare tło sygnalizuje, że
% można je pominąć przy powtórce do egzaminu. Zamiast oznaczać każdy dowód
% ręcznie, w wersji z wyróżnieniami opakowujemy CAŁE środowisko proof w szare
% pudełko — z wyjątkiem dowodów leżących już w ramce obowiazkowy (te są
% istotne i mają tło cream). Flaga \c@inobow śledzi, czy jesteśmy w obowiazkowy.
% W wersji zwykłej \tcolorboxenvironment nie jest wołane → proof bez zmian.
% inobow śledzi „jesteśmy w materiale obowiązkowym”. Ustawiamy go na środowisku
% obowiazkowy (a nie obowiazkowybox), żeby działał też bez wyróżnień — wersja
% z pominięciem dowodów potrzebuje tej informacji niezależnie od AZ_HIGHLIGHT.
\newbool{inobow}
\AtBeginEnvironment{obowiazkowy}{\booltrue{inobow}}
\AtEndEnvironment{obowiazkowy}{\boolfalse{inobow}}
\tcbset{
	pomijalnystyle/.style={%
		breakable, enhanced,
		colback=pomijalnybg,
		colframe=pomijalnyline,
		coltext=pomijalnytext,
		boxrule=0.5pt,
		arc=2pt, outer arc=2pt,
		left=10pt, right=10pt, top=8pt, bottom=8pt,
		boxsep=0pt,
		pad at break*=6pt,
		before skip=10pt plus 2pt, after skip=10pt plus 2pt,
	},
	% dowód w ramce obowiazkowy: bez szarego pudełka (zostaje na tle cream)
	pomijalnypusty/.style={blanker, before skip=\topsep, after skip=\topsep},
	% Rozstrzygamy styl jako KOD (a nie listę kluczy) w chwili czytania opcji
	% pudełka — wtedy \ifbool widzi już ustawione inobow (otwarte obowiazkowybox),
	% a problem dzielenia listy po przecinkach znika.
	proofbox/.code={\ifbool{inobow}{\tcbset{pomijalnypusty}}{\tcbset{pomijalnystyle}}},
}
% AZ_SKIP wygrywa: nieobowiązkowe dowody znikają zupełnie (przechwytujemy treść
% przez \BODY i nie składamy jej). Obowiązkowe składamy normalnie oryginalnym
% środowiskiem proof. Inaczej — w trybie wyróżnień — szare pudełka jak wyżej.
\ifbool{azskip}{%
	\let\azoldproof\proof
	\let\azoldendproof\endproof
	\RenewEnviron{proof}[1][\proofname]{%
		\ifbool{inobow}{\azoldproof[#1]\BODY\azoldendproof}{}%
	}%
}{\ifbool{azhighlight}{\tcolorboxenvironment{proof}{proofbox}}{}}

\titleformat{\section}{\large\bfseries\color{accent}}{}{0em}{}[\titlerule]
\titleformat{\subsection}{\normalsize\bfseries}{}{0em}{}
\titleformat{\subsubsection}{\normalsize\itshape}{}{0em}{}

\pagestyle{fancy}
\fancyhf{}

\rhead{\small \rightmark}
\lhead{\small \textit{Algorytmy Zaawansowane}}
\cfoot{\thepage}
\renewcommand{\headrulewidth}{0.4pt}

\hypersetup{colorlinks=true, linkcolor=accent, urlcolor=accent}

% ── KOMENDY ──────────────────────────────────────────────────────
% Wspólny nagłówek twierdzeń: gdy podano nazwę (argument [opcjonalny]),
% wypisz ją w nawiasie i przełam wiersz przed treścią; bez nazwy nagłówek
% pozostaje w jednej linii (run-in). \azThmHead{name}{number}{note}.
\newcommand{\azThmHead}[3]{%
	\thmname{#1}\thmnumber{ #2}%
	\def\azNote{#3}%
	\ifx\azNote\empty. \else\ (#3).\newline\fi}

% Większy odstęp przed/po + wspólny nagłówek z przełamaniem dla nazwanych.
% Wariant „definition” (treść prosta) i „plain” (treść kursywą).
\newtheoremstyle{azdefinition}%
{14pt plus 3pt minus 2pt}{12pt plus 3pt minus 2pt}%  odstęp przed / po
{\normalfont}%                                        czcionka treści
{0pt}%                                                wcięcie
{\bfseries}%                                          czcionka nagłówka
{}{0pt}%                                              interpunkcja i odstęp obsłużone w \azThmHead (headsep = 0)
{\azThmHead{#1}{#2}{#3}}%                          nagłówek
\newtheoremstyle{azplain}%
{14pt plus 3pt minus 2pt}{12pt plus 3pt minus 2pt}%
{\itshape}%
{0pt}%
{\bfseries}%
{}{0pt}%
{\azThmHead{#1}{#2}{#3}}%

\theoremstyle{azdefinition}
\newtheorem{definition}{Definicja}[section]
\newtheorem{example}[definition]{Przykład}
\newtheorem{problem}[definition]{Problem}

\theoremstyle{azplain}
\newtheorem{properties}[definition]{Własności}
\newtheorem{theorem}[definition]{Twierdzenie}
\newtheorem{conclusion}[definition]{Wniosek}
\newtheorem{lemma}[definition]{Lemat}

\theoremstyle{remark}
\newtheorem*{remark}{Uwaga}

% Nie pozostawiaj osamotnionego nagłówka twierdzenia na końcu strony:
% zarezerwuj kilka wierszy, w razie potrzeby przełam stronę przed środowiskiem.
\usepackage{needspace}
\newcommand{\azNeedSpace}{\Needspace*{4\baselineskip}}
\AtBeginEnvironment{definition}{\azNeedSpace}
\AtBeginEnvironment{example}{\azNeedSpace}
\AtBeginEnvironment{problem}{\azNeedSpace}
\AtBeginEnvironment{properties}{\azNeedSpace}
\AtBeginEnvironment{theorem}{\azNeedSpace}
\AtBeginEnvironment{conclusion}{\azNeedSpace}
\AtBeginEnvironment{lemma}{\azNeedSpace}

% Dygresja: wolniejsze rozwinięcie kroku dla czytelnika, który gubi się w rachunkach.
% Wyszarzona treść sygnalizuje materiał opcjonalny (można pominąć przy powtórce).
\newtheoremstyle{dygresjastyle}%
{\topsep}{\topsep}%                       odstęp przed / po
{\color{dygresjagray}}%                    czcionka treści (wyszarzona)
{0pt}%                                     wcięcie
{\color{dygresjagray}\itshape}%            czcionka nagłówka
{.}{.5em}{}%                               interpunkcja, odstęp po nagłówku, spec.
\theoremstyle{dygresjastyle}
\newtheorem*{dygresja}{Dygresja}


% ── KOMENDY ──────────────────────────────────────────────────────
% Podkreślenie terminów, które łamie się między wierszami (zwykłe \underline nie łamie).
\newcommand{\uemph}[1]{\uline{#1}}

\newcommand{\N}[0]{\mathbb{N}}
\newcommand{\Np}[0]{\mathbb{N_+}}
\newcommand{\R}[0]{\mathbb{R}}
\newcommand{\Rp}[0]{\mathbb{R_+}}

\newcommand{\BO}[0]{\mathcal{O}}
\newcommand{\BM}[0]{\mathcal{\Omega}}
\newcommand{\BT}[0]{\mathcal{\Theta}}

% Podwójne przekreślenie w „zwykłym” kierunku (//) — pakiet cancel daje tylko
% \cancel (/), \bcancel (\) i \xcancel (X). Oznacza skrócenie z wyrazem
% schowanym w wielokropku (drugi partner nie jest pokazany).
\newcommand{\dcancel}[1]{{\mathrlap{\cancel{#1}}\mkern8mu\cancel{\phantom{#1}}}}
% To samo w kierunku wstecznym (\\) — dla kompletu.
\newcommand{\dbcancel}[1]{{\mathrlap{\bcancel{#1}}\mkern8mu\bcancel{\phantom{#1}}}}

% \binom z XITS Math nie skaluje nawiasów wbudowanych w \genfrac — dolny
% argument wystaje pod nawias. Owijamy bezkreskowy \genfrac w prawdziwe
% \left(...\right), których wysokość dopasowuje się do dwuwierszowej zawartości.
\renewcommand{\binom}[2]{\left(\genfrac{}{}{0pt}{}{#1}{#2}\right)}

% ── KOMENDY ──────────────────────────────────────────────────────
% change default thickness of brackets to .6pt
\MHInternalSyntaxOn
\xpatchcmd\upbracketfill
{\sbox\z@{$\braceld$}\edef\l_MT_bracketheight_fdim{\the\ht\z@}}
{\edef\l_MT_bracketheight_fdim{.6pt}}
{}{\fail}

\xpatchcmd\downbracketfill
{\sbox\z@{$\braceld$}\edef\l_MT_bracketheight_fdim{\the\ht\z@}}
{\edef\l_MT_bracketheight_fdim{.6pt}}
{}{\fail}
\MHInternalSyntaxOff

% ── KOMENDY ──────────────────────────────────────────────────────

\algnewcommand{\To}[0]{\textbf{to} }
\algnewcommand{\Assign}[0]{$\;\,\gets\,\;$}
\algnewcommand{\Length}[1]{\textit{length}(#1)}
\algnewcommand{\Swap}[2]{\textit{swap}(#1, #2)}
\algnewcommand{\Size}[1]{\textit{size}(#1)}

\makeatletter
\newcommand{\FunctionNN}[2]{%
	\let\OldAlgLineNum\alglinenumber%
	\renewcommand{\alglinenumber}[1]{}%
	\Function{#1}{#2}%
	\let\alglinenumber\OldAlgLineNum%
	\addtocounter{ALG@line}{-1}% 
}
\makeatother

\makeatletter
\newenvironment{alg}[3][]{% #1=optional tcolorbox options, #2=caption, #3=label
	\if@noskipsec \leavevmode \fi% flush a pending run-in heading (\paragraph) so it stays above the box
	\par
	\addvspace{\bigskipamount}% breathing room before the box
	\refstepcounter{algorithm}\label{#3}% step the float counter + set ref/hyperref target
	\addcontentsline{\csname ext@algorithm\endcsname}{algorithm}% list-of-algorithms entry
	{\protect\numberline{\thealgorithm}{\ignorespaces #2}}%
	\small%
	\begin{tcolorbox}[
			breakable,
			sharp corners,
			boxrule=1pt,
			colback=white,
			colframe=black,
			colbacktitle=accent,
			coltitle=white,
			fonttitle=\bfseries,
			title={\csname fnum@algorithm\endcsname:\ \mdseries #2},% "Algorytm N: caption"
			left=2mm,
			right=2mm,
			top=2mm,
			bottom=2mm,
			#1%                    pass through any caller-supplied options
		]%
		\begin{algorithmic}[1]%
			}{%
		\end{algorithmic}%
	\end{tcolorbox}%
	\par\addvspace{\bigskipamount}% breathing room after the box
}
\makeatother

% ── RYSUNKI I TABELE BEZ PŁYWANIA ────────────────────────────────
% Zamiast pływających środowisk figure/table (i tak używanych prawie zawsze
% z [H], co przekreśla sens pływania) — wstawiamy treść dokładnie tam, gdzie
% jest napisana. Dzięki temu rysunek/tabela może też znaleźć się WEWNĄTRZ
% wyróżnienia obowiazkowy (pływak by z niego uciekł). \captionsetup{type=...}
% sprawia, że zwykłe \caption / \caption* numeruje się i etykietuje (\label,
% spis) tak samo jak w pływaku — tyle że bez pływania.
\newenvironment{azfigure}{%
	\par\addvspace{\bigskipamount}%
	\begingroup
	\captionsetup{type=figure}%
	\centering
}{%
	\par
	\endgroup
	\addvspace{\bigskipamount}%
}
\newenvironment{aztable}{%
	\par\addvspace{\bigskipamount}%
	\begingroup
	\captionsetup{type=table}%
	\centering
}{%
	\par
	\endgroup
	\addvspace{\bigskipamount}%
}
% ─────────────────────────────────────────────────────────────────
